\documentclass[11pt,a4paper]{article}
\usepackage[margin=0.75in]{geometry}
\usepackage{booktabs}
\usepackage{tabularx}
\usepackage{amsmath,amssymb}
\usepackage{xcolor}
\usepackage{hyperref}

\hypersetup{colorlinks=true, linkcolor=blue!70!black, urlcolor=blue!70!black}

\definecolor{boxgray}{gray}{0.94}
\definecolor{tagblue}{RGB}{28, 90, 160}

\newcommand{\reporttag}[1]{%
  \vspace{0.4em}\noindent\colorbox{tagblue!10}{\textbf{\textcolor{tagblue}{[Copy to Report: #1]}}}\vspace{0.2em}%
}

\title{\vspace{-1.5em}\Large\textbf{CSC311 Final Project: Modeling \& Technical Summary}\vspace{-0.5em}}
\author{\textbf{Group 63654}}
\date{\small November 2026 \quad|\quad Branch: \texttt{prediction}}

\begin{document}
\maketitle
\vspace{-1.5em}

This technical summary documents the preprocessing pipeline, hyperparameter tuning experiments, evaluation benchmarks, and standalone deployment verification for the CSC311 sports classification project.

% ==============================================================================
\section{Pipeline Implementation \& Module Summary}
% ==============================================================================
The project implementation consists of three self-contained modules:
\begin{enumerate}
    \item \textbf{Data Cleaning \& Feature Engineering (\texttt{src/data/data\_cleaning.py})}:
    \begin{itemize}
        \item \textbf{TF-IDF Vocabulary}: Set \texttt{max\_features=25} per free-text column, capturing informative sport keywords while eliminating single-occurrence spelling variations and noise.
        \item \textbf{Numeric Standardization}: Computed Z-score scaling strictly on training respondents for ticket price ($\log(p+1)$) and difficulty rating. This eliminates feature scale discrepancies and guarantees stable, convex optimization.
        \item \textbf{Categorical Encoding}: Transformed multi-select responses (season, venue, music) into binary indicator matrices and mapped Likert items to 1--5 ordinal values.
    \end{itemize}
    \item \textbf{Hyperparameter Grid Search \& Evaluation (\texttt{src/models/train\_and\_tune.py})}:
    \begin{itemize}
        \item Implemented systematic tuning across Softmax Regression, Random Forest, and Multinomial Naive Bayes.
        \item Automatically computes and reports Train Accuracy, Validation Accuracy, Macro-F1, Precision, Recall, and Confusion Matrices.
    \end{itemize}
    \item \textbf{Standalone Inference Engine (\texttt{src/prediction/pred.py})}:
    \begin{itemize}
        \item Exposes \texttt{predict\_all(csv\_path)} taking unseen test CSVs and returning predictions as a Python string list.
        \item \textbf{Zero External Dependencies}: Implemented strictly using standard Python, NumPy, and Pandas (no \texttt{sklearn}, \texttt{torch}, or \texttt{scipy}).
        \item \textbf{Pre-trained Artifacts}: Model weights ($W \in \mathbb{R}^{3 \times 133}, b \in \mathbb{R}^3$) and vocabulary mappings are serialized into \texttt{model\_bundle.json} (23 KB, well below the 10 MB limit).
    \end{itemize}
\end{enumerate}

% ==============================================================================
\section{Hyperparameter Tuning Decisions: Why These Exact Numbers?}
% ==============================================================================
\reporttag{Section 4: Methodology --- Hyperparameter Search}

Following course guidelines, all three candidate model families were systematically tuned over discrete grids before model selection:

\subsection{Model 1: Softmax Regression --- L2 Penalty (C)}
Inverse regularization parameter $C = 1/\lambda$ was tested across $C \in [0.01, 0.1, 0.5, 1.0, 5.0, 10.0]$:
\begin{center}
\small
\begin{tabularx}{\textwidth}{ccclX}
\toprule
$C$ & \textbf{Train Acc} & \textbf{Val Acc} & \textbf{Diagnostic State} & \textbf{Why Chosen or Rejected?} \\
\midrule
0.01 & 89.00\% & 88.37\% & Underfitting & L2 penalty too aggressive; heavily shrinks useful weights. \\
0.10 & 94.09\% & 93.80\% & Moderate & Reasonable boundary, but under-utilizes subtle text signals. \\
0.50 & 94.75\% & 94.57\% & Near-optimal & Good balance, slightly lower training separation. \\
\textbf{1.00} & \textbf{94.91\%} & \textbf{94.57\%} & \textbf{Optimal} & \textbf{Selected}: Peak validation accuracy with minimal gap to train. \\
5.00 & 95.73\% & 94.57\% & Slight Overfit & Slight train gain without val improvement. \\
10.00 & 96.22\% & 94.57\% & Overfitting & Regularization too weak; risks fitting outlier respondents. \\
\bottomrule
\end{tabularx}
\end{center}

\subsection{Model 2: Random Forest --- Maximum Tree Depth (max\_depth)}
Evaluated tree depth limit $\text{max\_depth} \in [5, 10, 15, \text{None}]$ with 100 estimators:
\begin{center}
\small
\begin{tabularx}{\textwidth}{ccclX}
\toprule
\textbf{max\_depth} & \textbf{Train Acc} & \textbf{Val Acc} & \textbf{Diagnostic State} & \textbf{Why Chosen or Rejected?} \\
\midrule
5 & 95.40\% & 93.80\% & Underfitting & Trees too shallow to capture non-linear joint interactions. \\
\textbf{10} & \textbf{97.87\%} & \textbf{95.35\%} & \textbf{Optimal} & \textbf{Selected}: Balances rich interaction modeling with tree pruning. \\
15 & 99.84\% & 95.35\% & Memorization & Train acc reaches 99.8\%; no validation gain over depth 10. \\
None & 100.0\% & 94.57\% & Severe Overfit & Pure leaves memorize training set; validation drops by 0.78\%. \\
\bottomrule
\end{tabularx}
\end{center}

\subsection{Model 3: Multinomial Naive Bayes --- Laplace Smoothing (Alpha)}
Evaluated additive smoothing $\alpha \in [0.1, 0.5, 1.0, 2.0]$ on free-text TF-IDF features:
\begin{center}
\small
\begin{tabularx}{\textwidth}{ccclX}
\toprule
$\alpha$ & \textbf{Train Acc} & \textbf{Val Acc} & \textbf{Diagnostic State} & \textbf{Why Chosen or Rejected?} \\
\midrule
\textbf{0.1} & \textbf{90.31\%} & \textbf{89.92\%} & \textbf{Optimal} & \textbf{Selected}: Prevents zero probabilities while keeping tokens sharp. \\
0.5 & 90.31\% & 89.92\% & Comparable & Similar validation accuracy. \\
1.0 & 89.82\% & 89.15\% & Over-smoothed & Standard Laplace smoothing slightly dampens rare keywords. \\
2.0 & 89.33\% & 89.15\% & Over-smoothed & Excessive smoothing dilutes likelihood contrast across sports. \\
\bottomrule
\end{tabularx}
\end{center}

% ==============================================================================
\section{Multi-Tier Testing \& Autograder Verification}
% ==============================================================================
The pipeline was verified through three evaluation stages:
\begin{enumerate}
    \item \textbf{Tier 1: Validation Set Evaluation ($N=129$)}:
    \begin{itemize}
        \item Softmax Regression ($C=1.0$): \textbf{94.57\%} Accuracy, \textbf{94.57\%} Macro-F1.
        \item Random Forest ($\text{depth}=10$): \textbf{95.35\%} Accuracy, \textbf{95.37\%} Macro-F1.
    \end{itemize}
    \item \textbf{Tier 2: Held-out Internal Test Set ($N=132$, strictly sequestered until final run)}:
    \begin{itemize}
        \item Both Softmax Regression and Random Forest scored exactly \textbf{89.39\%} accuracy.
        \item Given the random guessing baseline of $33.3\%$ and natural survey noise, $89.4\%$ confirms strong, un-overfitted generalization on completely unseen respondents.
    \end{itemize}
    \item \textbf{Tier 3: Autograder Robustness \& Stress Tests on \texttt{pred.py}}:
    \begin{itemize}
        \item \textbf{Missing/Null Edge Case}: Tested on an all-NaN row $\rightarrow$ handles median imputation safely.
        \item \textbf{Dirty String Edge Case}: Tested on currency strings (\texttt{"\$5,000 dollars"}, \texttt{"free"}, \texttt{"1k"}) $\rightarrow$ correct numeric parsing.
        \item \textbf{Runtime \& Footprint}: 870 rows processed in \textbf{0.9 seconds} ($<0.001$s per row, well under the 60 rows/min requirement). Total bundle size is \textbf{28 KB}.
    \end{itemize}
\end{enumerate}

% ==============================================================================
\section{Results Table \& Error Analysis}
% ==============================================================================
\reporttag{Section 5: Results --- Comparison Table \& Confusion Matrix}

\subsection{Summary Table (Direct Copy into Paper)}
\begin{center}
\small
\begin{tabular*}{\textwidth}{@{\extracolsep{\fill}}lccccc@{}}
\toprule
\textbf{Model Family} & \textbf{Train Acc} & \textbf{Val Acc} & \textbf{Val Macro-F1} & \textbf{Test Acc} & \textbf{Test Macro-F1} \\
\midrule
\textbf{Softmax Regression} ($C=1.0$) & 94.91\% & 94.57\% & 94.57\% & \textbf{89.39\%} & \textbf{89.42\%} \\
\textbf{Random Forest} ($\text{depth}=10$) & 97.87\% & 95.35\% & 95.37\% & \textbf{89.39\%} & \textbf{89.39\%} \\
\textbf{Naive Bayes} ($\alpha=0.1$, text) & 90.31\% & 89.92\% & 89.94\% & 84.85\% & 84.87\% \\
\bottomrule
\end{tabular*}
\end{center}

\subsection{Confusion Matrix (Validation Set, $N=129$)}
\begin{center}
\small
\begin{tabular}{lccc}
\toprule
& \textbf{Pred: Soccer} & \textbf{Pred: Baseball} & \textbf{Pred: Basketball} \\
\midrule
\textbf{True: Soccer} & \textbf{40} & 1 & 2 \\
\textbf{True: Baseball} & 2 & \textbf{41} & 0 \\
\textbf{True: Basketball} & 1 & 1 & \textbf{41} \\
\bottomrule
\end{tabular}
\end{center}

\subsection{Qualitative Error Breakdown}
\begin{itemize}
    \item \textbf{Soccer vs.\ Baseball Overlap}: The majority of misclassifications occur between Soccer and Baseball (2 in Val, 6 in Test). Both sports share outdoor fields (\textit{parks, community recreation}), Spring/Summer seasonality, and similar ticket willingness (\$15--\$40).
    \item \textbf{Basketball Distinctiveness}: Basketball has the lowest error rate across all splits. Strong discriminative features include indoor facilities, winter seasonality, and distinct vocabulary (\textit{dunk, sneakers, court, nba}).
\end{itemize}

% ==============================================================================
\section{Report Final Requirement: Projected Test Set Accuracy}
% ==============================================================================
\reporttag{Section 5: Final Single-Number Projection}

\noindent\textbf{Projected Course Test Accuracy}: \textbf{89.4\%} \\
\textbf{Theoretical Justification}: Justified directly by the sequestered internal test set result (89.39\%). Since hyperparameter tuning was conducted exclusively on the validation set, the internal test score is an unbiased estimate of generalizability on unseen respondents.

\end{document}
